\documentclass[12pt,a4paper]{article}
        \makeatletter
        \def\input@path{{../}}
        \makeatother
        \usepackage{almeria-fichas}

        \FichaMeta{Sesion 06}{Figuras planas en la ciudad}{Bloque 2 · Geometria urbana}{Clasificar triangulos, cuadrilateros y circunferencias presentes en modelos urbanos.}{Conceptual}{Recordar, Comprender, Analizar}

        \begin{document}

        \FichaCabecera

        \begin{materialesbox}
        \begin{itemize}
    \item Ficha impresa
    \item Lapiz
    \item Regla
    \item Lapis de colores
\end{itemize}
        \end{materialesbox}

        \begin{pasosbox}
        \begin{enumerate}
    \item Observa la forma general antes de nombrarla.
    \item Cuenta lados y vertices si lo necesitas.
    \item Escribe siempre una propiedad para justificar.
    \item Haz un dibujo sencillo cuando tengas dudas.
\end{enumerate}
        \end{pasosbox}

        \begin{recordatoriobox}
        \begin{itemize}
    \item Los triangulos tienen 3 lados.
    \item Los cuadrilateros tienen 4 lados.
    \item Una circunferencia es la linea; el circulo es la superficie interior.
\end{itemize}
        \end{recordatoriobox}

        \begin{apoyobox}
        \begin{itemize}
    \item Subraya la palabra que te da la pista: lados, angulos o lados iguales.
    \item Puedes clasificar primero y justificar despues.
    \item Si dos figuras se parecen, compara el numero de lados.
\end{itemize}
        \end{apoyobox}

        \BloomSection{recordar}

\BloomExercise{recordar}{1}{Nombres de figuras}{Escribe el nombre de cinco figuras planas que conozcas.}{5}

\BloomExercise{recordar}{2}{Propiedades basicas}{Completa: un triangulo tiene \underline{\hspace{2cm}} lados y un cuadrilatero tiene \underline{\hspace{2cm}} lados.}{4}

\BloomSection{comprender}

\BloomExercise{comprender}{1}{Justifica la clasificacion}{Explica por que un cuadrado es tambien un cuadrilatero y por que un triangulo nunca puede tener cuatro vertices.}{6}

\BloomExercise{comprender}{2}{Circulo o circunferencia}{Escribe una frase para distinguir \textbf{circulo} y \textbf{circunferencia}.}{5}

\BloomSection{analizar}

\BloomExercise{analizar}{1}{Comparacion de cuadrilateros}{Analiza en que se parecen y en que se diferencian un rectangulo y un rombo.}{7}

\BloomExercise{analizar}{2}{Eleccion de forma}{Una plaza necesita una zona central visible desde todos sus bordes. ¿Seria mejor una forma circular o rectangular? Analiza ventajas y limites.}{7}

        \begin{evidenciabox}
        \begin{itemize}
    \item Reconocer figuras planas en contextos urbanos.
    \item Distinguir triangulos y cuadrilateros por propiedades.
    \item Explicar por que una figura pertenece a una categoria.
\end{itemize}
        \end{evidenciabox}

        \end{document}
